\documentclass[10pt, twocolumn, a4paper]{article}

\usepackage{cite}
\usepackage{amsmath,amssymb,amsfonts}
\usepackage{algorithmic}
\usepackage{graphicx}
\usepackage{textcomp}
\usepackage{xcolor}
\usepackage{url}
\usepackage{hyperref}
\usepackage[utf8]{inputenc}
\usepackage[portuguese]{babel}
\usepackage{authblk} % Para formatação de múltiplos autores
\usepackage{geometry}
\geometry{a4paper, total={170mm,257mm}, left=20mm, top=20mm}

\title{Título do Seu Artigo Aqui: Uma Abordagem Abrangente}

\author[1]{Primeiro A. Autor}
\author[2]{Segundo B. Autor}
\author[2]{Terceiro C. Autor}

\affil[1]{Instituto de Exemplo, Cidade, Estado, CEP (e-mail: autor@exemplo.com)}
\affil[2]{Departamento de Ciência da Computação, Universidade de Exemplo, Cidade, Estado, CEP}

\date{} % Deixe vazio para omitir a data

\begin{document}

\maketitle

\begin{abstract}
Este é o resumo do seu artigo. Ele deve sumarizar de forma concisa o problema que você está resolvendo, a metodologia proposta e os principais resultados ou contribuições. O resumo é geralmente um único parágrafo com cerca de 150 a 250 palavras. Não cite referências no resumo.
\end{abstract}

\vspace{1em}
\noindent\textbf{Palavras-chave:} Deep Learning, Detecção de Incêndio, \LaTeX, Sensoriamento Remoto, Detecção de Fumaça.
\vspace{1em}

\section{Introdução}
Esta é a seção de introdução. Aqui você apresenta o contexto, a motivação, a definição do problema e resume suas contribuições. 

Grandes editoras esperam uma introdução forte que estabeleça o contexto do trabalho e declare claramente por que é importante resolver o problema abordado. 

\subsection{Motivação}
Por que este problema é importante? Por exemplo, a detecção de incêndios por sensoriamento remoto é crucial para a proteção ambiental e o gerenciamento de desastres.

\subsection{Contribuições}
As principais contribuições deste artigo são resumidas a seguir:
\begin{itemize}
    \item Apresentamos um novo conjunto de dados (dataset) para detecção de chamas abertas e fumaça.
    \item Propomos uma arquitetura de aprendizado profundo que alcança o estado da arte em precisão.
    \item Realizamos avaliações extensivas sobre a abordagem proposta.
\end{itemize}

\section{Trabalhos Relacionados}
Discuta a literatura relevante aqui. Agrupe os trabalhos relacionados de forma lógica, por exemplo, pelas técnicas utilizadas ou pelos domínios de aplicação. Você pode usar o comando \verb|\cite{}| para referenciar artigos como a pesquisa do dataset FASDD \cite{exemplo_ref}.

Caso precise citar livros-texto fundamentais, o processo é o mesmo. Por exemplo, ao explicar os conceitos base de redes neurais, é muito comum citar o livro Deep Learning de Goodfellow et al. \cite{goodfellow2016deep}. Para manuais técnicos, você poderia citar o LaTeX Companion \cite{latex_companion}. Se você quiser citar várias referências de uma só vez, o LaTeX permite agrupá-las \cite{exemplo_ref, goodfellow2016deep, latex_companion}.

\section{Metodologia}
Detalhe seu método proposto, sua estrutura (framework) ou o conjunto de dados criado aqui.

\subsection{Coleta e Pré-processamento do Dataset}
Descreva as etapas de coleta, limpeza e pré-processamento de dados aqui.

\subsection{Arquitetura Proposta}
Descreva a formulação matemática e arquitetônica de seu modelo. 

\begin{equation}
    \mathcal{L}_{total} = \alpha \mathcal{L}_{cls} + \beta \mathcal{L}_{reg} \label{eq:loss}
\end{equation}

Como mostrado na Eq. \ref{eq:loss}, você pode facilmente fazer referências cruzadas de equações no LaTeX.

\section{Experimentos e Resultados}
Descreva a configuração experimental, os métodos de base (baselines) e apresente seus resultados quantitativos e qualitativos.

\subsection{Configuração Experimental}
Detalhes sobre hardware, software, conjuntos de dados e hiperparâmetros.

\subsection{Resultados}
Você pode incluir figuras e tabelas para mostrar seus resultados. 

\begin{figure}[htbp]
\centerline{\includegraphics[width=0.45\textwidth]{example-image}}
\caption{Exemplo de legenda de figura. Certifique-se de que as figuras estejam legíveis.}
\label{fig:arquitetura}
\end{figure}

\begin{table}[htbp]
\caption{Comparação de Desempenho dos Modelos}
\begin{center}
\begin{tabular}{|c|c|c|c|}
\hline
\textbf{Modelo} & \textbf{Precisão} & \textbf{Revocação (Recall)} & \textbf{F1-Score} \\
\hline
Baseline A & 0,85 & 0,82 & 0,83 \\
\hline
Baseline B & 0,88 & 0,86 & 0,87 \\
\hline
\textbf{Proposto} & \textbf{0,94} & \textbf{0,92} & \textbf{0,93} \\
\hline
\end{tabular}
\label{tab:resultados}
\end{center}
\end{table}

\section{Discussão}
Discuta as implicações de seus resultados, as limitações de seu estudo e os potenciais trabalhos futuros. Por que seu modelo teve um desempenho melhor? Em que cenários ele falha?

\section{Conclusão}
Conclua seu artigo resumindo as principais descobertas e a sua relevância para a área de pesquisa.

\begin{thebibliography}{99}

\bibitem{exemplo_ref}
M. Wang, P. Yue, L. Jiang, D. Yu, T. Tuo, e J. Li, ``An open flame and smoke detection dataset for deep learning in remote sensing based fire detection,'' \textit{Geo-spatial Information Science}, vol. 28, no. 2, pp. 511--526, 2025.

\bibitem{goodfellow2016deep}
I. Goodfellow, Y. Bengio, e A. Courville, \textit{Deep Learning}. MIT press, 2016.

\bibitem{latex_companion}
F. Mittelbach, M. Goossens, J. Braams, D. Carlisle, e C. Rowley, \textit{The \LaTeX{} Companion}, 2nd ed. Addison-Wesley Professional, 2004.

\end{thebibliography} 

\end{document}
